% Analogías de las lecciones de Comprensión lectora (libros 3 y 4). Salen solas después del reto de cada lección.
% --- Libro 3 ---
\analogiade{pu-el-1}{El tráiler de la película}{Nadie entra a cine sin saber si la película es de terror o de risa:
el tráiler lo anuncia. La prelectura es el tráiler del texto. Treinta segundos mirando el título, el comienzo de cada
párrafo y el final te dicen qué clase de historia viene, y así nada te toma por sorpresa.}
\analogiade{pu-el-2}{El mapa del centro comercial}{El mapa de un centro comercial no describe cada tienda: dice
«comidas aquí, cines allá, salida al fondo». El mapa del texto hace lo mismo con los párrafos: «aquí presenta, aquí
da un ejemplo, aquí concluye». Con ese mapa, cuando una pregunta te manda a buscar algo, sabes a qué piso ir.}
\analogiade{pu-el-3}{La lista de invitados}{En una pregunta EXCEPTO eres el portero con la lista: revisas uno por uno
y dejas pasar a los que están anotados en el texto. El que no aparece en la lista es justo el que buscas. No se trata
de encontrar al invitado más importante, sino al colado.}
\analogiade{pu-tt-1}{Las herramientas de la caja}{Un martillo, un destornillador y una llave sirven para cosas
distintas. Los tipos de texto son herramientas: uno cuenta, otro describe, otro explica, otro convence, otro da
instrucciones. Antes de leer, pregúntate para qué fue hecha esa herramienta, y sabrás cómo usarla.}
\analogiade{pu-tt-2}{Los géneros musicales}{Con los primeros compases sabes si suena salsa, rock o vallenato, y ya
esperas cierto ritmo y ciertos instrumentos. Con los géneros de texto pasa igual: reconocer que es una reseña, una
crónica o un editorial te dice qué partes va a tener y qué te van a preguntar.}
\analogiade{pu-tt-3}{El esqueleto}{Dos personas muy distintas tienen el mismo esqueleto por dentro. Dos textos muy
distintos pueden tener el mismo esqueleto: problema y solución, causa y efecto, antes y después. Ver el esqueleto bajo
la ropa de las palabras es entender cómo está armado el texto.}
\analogiade{pu-ie-1}{Buscar un contacto en el celular}{Para llamar a alguien no lees toda la lista de contactos:
escribes las primeras letras y vas directo. Para ubicar un dato en el texto, igual: eliges una palabra ancla de la
pregunta, la buscas con la vista y lees con cuidado solo esa zona.}
\analogiade{pu-ie-2}{El mismo regalo con otro papel}{Una paráfrasis es el mismo regalo envuelto en otro papel. La
respuesta correcta casi nunca trae el papel original (las mismas palabras del texto); trae el mismo regalo. Y las
opciones tramposas hacen lo contrario: usan el papel del texto para envolver otra cosa.}
\analogiade{pu-ie-3}{El testigo en el juicio}{Un buen testigo dice «eso lo vi», «eso no pasó» o «eso no lo sé». La
tercera respuesta es tan importante como las otras. Frente a cada opción, sé ese testigo: ¿el texto lo dice, lo niega
o simplemente no habla de eso?}
\analogiade{pu-ie-4}{El chat grupal}{En un chat de grupo no basta saber \emph{qué} se dijo: importa \emph{quién} lo
dijo. Atribuirle a tu mamá el mensaje de tu tío cambia toda la historia. Los textos con varias voces son chats: marca
quién habla en cada parte antes de responder.}
\analogiade{pu-vo-1}{La ropa según la ocasión}{No te vistes igual para un grado que para un partido. Las palabras
tampoco: «acoger» significa una cosa junto a «un viajero» y otra junto a «una propuesta». El sinónimo correcto es el
que le queda bien a \emph{esa} ocasión: pruébaselo a la frase y mira si le luce.}
\analogiade{pu-vo-2}{La llave que abre varias puertas}{Hay palabras que son llaves maestras: «banco», «sierra»,
«planta». Abren varias puertas, pero en cada frase solo una. Las palabras vecinas te dicen frente a qué puerta estás:
junto a «dinero», el banco es uno; junto a «parque», es otro.}
\analogiade{pu-vo-3}{El apodo}{Si a un amigo muy alto le dicen «el Enano» o a un calvo «el Peludo», nadie lo toma al pie
de la letra: todos saben a quién se refieren. Las frases hechas son apodos de ideas. «Tirar la toalla» no habla de
toallas; habla de rendirse. Quien la lee literalmente no entendió el apodo.}
\analogiade{pu-vo-4}{Las fichas de armar}{Las palabras largas están hechas de fichas que se repiten: \emph{tele}
(lejos), \emph{bio} (vida), \emph{logía} (estudio). Si conoces las fichas, puedes armar y desarmar palabras que nunca
habías visto, como quien reconoce las piezas en una construcción nueva.}
\analogiade{pu-rf-1}{El dedo que señala}{Un pronombre es un dedo que señala: «ella», «eso», «aquel». Solo no dice
nada; hay que mirar hacia dónde apunta. Casi siempre apunta hacia atrás, a alguien que ya fue nombrado. Si el dedo
puede señalar a dos, prueba cada uno y quédate con el que da sentido.}
\analogiade{pu-rf-2}{Lo que no hace falta repetir}{Si te preguntan «¿quieres café?» respondes «sí, con leche», no «sí,
quiero café con leche». Lo demás se sobreentiende. Los textos hacen eso todo el tiempo: omiten lo que ya dijeron. Leer
bien es reponer en la mente la parte que el autor se ahorró.}
\analogiade{pu-cn-1}{Las señales de tránsito}{Los conectores son las señales de la carretera: «además» es siga
derecho; «sin embargo», curva cerrada; «por lo tanto», llegada; «por ejemplo», mirador. Quien lee las señales sabe
qué viene antes de leerlo, y no se sale en la curva.}
\analogiade{pu-cn-2}{Los jugadores del equipo}{En un equipo, cada jugador tiene un puesto: uno ataja, otro arma, otro
remata. En un texto, cada párrafo tiene su puesto: uno presenta, otro da un ejemplo, otro responde a una objeción,
otro cierra. La pregunta «¿qué función cumple?» es «¿en qué puesto juega?».}
% --- Libro 4 ---
\analogiade{pu-in-1}{Las huellas en el barro}{Si ves huellas de botas mojadas que entran a la casa, sabes que alguien llegó y que venía de un lugar mojado, aunque nadie te lo haya dicho. Inferir es eso: leer huellas. Pero un buen rastreador no
inventa: solo concluye lo que las huellas permiten.}
\analogiade{pu-in-2}{Los cimientos de la casa}{No se ven, pero sin ellos la casa no se sostiene. Los presupuestos son
los cimientos de una frase: «volvió a llegar tarde» se apoya en que ya había llegado tarde antes. Para descubrirlos,
pregunta: ¿qué tiene que ser cierto para que esta frase tenga sentido?}
\analogiade{pu-in-3}{La fila de fichas de dominó}{Empujas la primera ficha y caen todas, una tras otra. En una cadena
de causas, cada hecho tumba al siguiente. Cuando preguntan «¿por qué cayó la última?», mira si quieren la ficha que
la tocó (la causa directa) o el dedo que empujó la primera (la causa original).}
\analogiade{pu-in-4}{El cine mudo}{En el cine mudo nadie dice «estoy triste»: lo ves en cómo camina el personaje, en
lo que hace con las manos. Los buenos textos son así. Para saber qué siente un personaje, no busques la palabra
«triste»: mira qué hace, qué calla y qué gestos repite.}
\analogiade{pu-ap-1}{La receta en otra cocina}{Si entendiste una receta de verdad, puedes hacerla en otra cocina, con
otra estufa y otras ollas. Aplicar una idea a un caso nuevo es eso: llevarte lo esencial (el principio) y reconocerlo
aunque cambien los ingredientes del ejemplo.}
\analogiade{pu-ap-2}{El puente}{Una analogía es un puente entre dos orillas que no se parecen por fuera pero funcionan
igual por dentro: el corazón y una bomba de agua, el maestro y el jardinero. Para cruzarlo, hay que nombrar la
relación exacta que une a la primera pareja y buscarla en la segunda.}
\analogiade{pu-fu-1}{Los invitados del programa}{El presentador de un programa invita gente: a unos para que apoyen
lo que él piensa, a otros para discutirles. No todo lo que dice un invitado es lo que piensa el presentador. En un
texto, el autor es el presentador: averigua para qué trajo a cada voz.}
\analogiade{pu-fu-2}{El vendedor y el mecánico}{Si el vendedor te dice que el carro está perfecto, le crees menos que
al mecánico independiente que lo revisó. No porque el vendedor mienta, sino porque gana si le crees. Con las fuentes
de un texto se hace la misma pregunta: ¿quién habla, qué sabe y qué gana?}
\analogiade{pu-ag-1}{La mesa y sus patas}{La tesis es la tabla de la mesa; los argumentos son las patas. Una mesa con
una sola pata débil se cae. Para entender un texto que quiere convencer, mira primero qué está sosteniendo (la tabla)
y luego cuenta y prueba las patas una por una.}
\analogiade{pu-ag-2}{El tono de voz}{«Qué bien», dicho con una sonrisa, felicita; dicho entre dientes, reclama. Las
mismas palabras, otro tono. En un texto escrito el tono se oye en los adjetivos, en lo que se exagera y en lo que se
calla. Leer el tono es escuchar la voz detrás de las letras.}
\analogiade{pu-ag-3}{El termómetro y el «hace frío»}{El termómetro marca 18 grados: eso es un hecho, cualquiera puede
comprobarlo. «Hace frío» es una opinión: para alguien de la costa sí, para alguien del páramo no. La prueba es
sencilla: ¿se puede verificar con un instrumento o un registro, o depende de quién lo diga?}
\analogiade{pu-lf-1}{El lente de aumento}{El lenguaje figurado es un lente: no cambia la realidad, la hace ver más
grande o desde otro ángulo. Decir «tengo un nudo en la garganta» no describe una cuerda: hace sentir la emoción. Ante
una figura, pregunta siempre qué se ve mejor gracias a ese lente.}
\analogiade{pu-lf-2}{El guiño}{La ironía es decir algo y picar el ojo. Si rompes un vaso y tu hermana dice «¡qué
delicadeza!», el guiño está en el choque entre sus palabras y el vaso roto. En un texto el guiño no se ve: se
descubre cuando lo que se dice no cuadra con lo que pasa.}
\analogiade{pu-dt-1}{Dos testigos del mismo partido}{Dos hinchas de equipos rivales cuentan el mismo partido. Coinciden
en el marcador y en que hubo un penalti; difieren en si fue justo. Leer dos textos es escucharlos a los dos: separar
en qué están de acuerdo (aunque no lo parezca) y en qué punto exacto se apartan.}
\analogiade{pu-dt-2}{Pensar como el otro}{Si conoces bien a tu abuela, sabes qué diría de tu nuevo corte de pelo sin
preguntarle. No es adivinar: es aplicar lo que sabes de cómo piensa. Con un autor es igual: su tesis y sus valores te
dicen cómo respondería a una idea que nunca comentó.}
